\documentclass{beamer}

\usepackage{amsmath}
\usepackage{amsfonts}
\usepackage{amssymb}
\usepackage{listings}
\usepackage{graphicx}
\usepackage{tikz}

% ===== KCPC clean theme =====
\usetheme{default}
\useinnertheme{rectangles}
\useoutertheme{default}

%\beamerdefaultoverlayspecification{<+->}

\setbeamertemplate{navigation symbols}{}

% --- Colours ---
\definecolor{kcpc-bg}{RGB}{250,250,250}
\definecolor{kcpc-text}{RGB}{20,20,20}
\definecolor{kcpc-muted}{RGB}{90,90,90}
\definecolor{kcpc-red}{RGB}{176,0,0}

\setbeamercolor{background canvas}{bg=kcpc-bg}
\setbeamercolor{normal text}{fg=kcpc-text}
\setbeamercolor{frametitle}{fg=kcpc-text,bg=kcpc-bg}
\setbeamercolor{title}{fg=kcpc-text}
\setbeamercolor{subtitle}{fg=kcpc-muted}
\setbeamercolor{structure}{fg=kcpc-red}

\setbeamerfont{title}{series=\bfseries,size=\Large}
\setbeamerfont{frametitle}{series=\bfseries,size=\large}

\setlength{\parskip}{0.4em}

\setbeamertemplate{itemize item}{→}
\setbeamertemplate{itemize subitem}{→}

\setbeamertemplate{headline}{
    \vspace{0.2cm}
    \hspace{0.2cm}
    \textcolor{kcpc-muted}{\small\insertsection}
    \hfill
}

\setbeamertemplate{footline}{
    \hspace{0.2cm}
    \textcolor{kcpc-muted}{\small KCPC}
    \hfill
    \vspace{0.3cm}
}

\setbeamertemplate{background}{
    \begin{tikzpicture}[remember picture,overlay]
        \node[
            opacity=0.035,
            at=(current page.center)
        ] {
            \includegraphics[width=0.55\paperwidth]{kcpc-logo.png}
        };
    \end{tikzpicture}
}

\AtBeginSection[]{
    \begin{frame}[plain]
        \vfill
        \centering
        {\Large\bfseries \insertsection}
        \vfill
    \end{frame}
}


\title{Graphs II}
\subtitle{Distance, Constraints, and Logical Structure}
\author{KCPC}
\date{\today}

\begin{document}

% ------------------------------------------------

\begin{frame}
    \titlepage
\end{frame}

% ------------------------------------------------

\begin{frame}{Overview}
    \tableofcontents
\end{frame}

% ------------------------------------------------

\begin{frame}{Attendance}
    \centering
    \includegraphics[width=0.5\paperwidth]{attendance.png}
\end{frame}

% ------------------------------------------------
\begin{frame}{Our Socials}

    \centering
    \begin{columns}[c]
    
        \column{0.33\textwidth}
        \centering
        
\includegraphics[width=\textwidth]{insta.JPG}
        \\
        Instagram
        
        \column{0.33\textwidth}
        \centering
        
\includegraphics[width=\textwidth]{discord.JPG}
        \\
        Discord
        
        \column{0.33\textwidth}
        \centering
        
\includegraphics[width=\textwidth]{whatsapp.JPG}
        \\
        WhatsApp
    
    \end{columns}

\end{frame}

% ------------------------------------------------

\section{Foundations - A Quick Graph Traversal Recap}

% ------------------------------------------------

\begin{frame}{Depth-First Search}

    \textbf{Purpose}
    
        Explore as far as possible along each branch before backtracking.
    
    \vspace{0.5em}
    
\end{frame}

\begin{frame}{Depth-First Search}

    \textbf{Algorithm (Iterative)}
    
    \textbf{1. Initialize}
        \begin{itemize}
            \item Create empty \texttt{stack}
            \item Mark all nodes unvisited
            \item Push start node
        \end{itemize}

\end{frame}

\begin{frame}{Depth-First Search}

    \textbf{2. While stack not empty}
        \begin{itemize}
            \item Pop node $u$
            \item If $u$ unvisited:
            \begin{itemize}
                \item Mark $u$ visited
                \item Push unvisited neighbours of $u$
            \end{itemize}
        \end{itemize}
    
\end{frame}

\begin{frame}{Depth-First Search}

    \textbf{3. Stop when stack empty}

\end{frame}

% ------------------

\begin{frame}{Group walkthrough example}

    % We assume the source (obviously) here is A, with us mapping
    % all possible routes until we find our target node.
    
    \includegraphics<1>[width=0.8\linewidth]{DFS1.jpg}
    \includegraphics<2>[width=0.8\linewidth]{DFS2.jpg}
    \includegraphics<3>[width=0.8\linewidth]{DFS3.jpg}
    \includegraphics<4>[width=0.8\linewidth]{DFS4.jpg}
    \includegraphics<5>[width=0.8\linewidth]{DFS5.jpg}
    \includegraphics<6>[width=0.8\linewidth]{DFS6.jpg}
    \includegraphics<7>[width=0.8\linewidth]{DFS7.jpg}

\end{frame}

% ------------------

\begin{frame}{Breadth-First Search (BFS)}

    \textbf{Purpose}
    
        Explore a graph level by level from a starting node.

\end{frame}

\begin{frame}{Breadth-First Search (BFS)}

    \vspace{0.5em}
    \textbf{Algorithm}
    
        \textbf{1. Initialize}
        \begin{itemize}
            \item Create empty \texttt{queue}
            \item Mark all nodes unvisited
            \item Mark start node visited
            \item Enqueue start node
        \end{itemize}
    
\end{frame}
    
\begin{frame}{Breadth-First Search (BFS)}
    
    \textbf{2. While queue not empty}
        \begin{itemize}
            \item Dequeue node $u$
            \item For each neighbour $v$ of $u$:
            \begin{itemize}
                \item If $v$ unvisited:
                \begin{itemize}
                    \item Mark $v$ visited
                    \item Enqueue $v$
                \end{itemize}
            \end{itemize}
        \end{itemize}

\end{frame}

\begin{frame}{Breadth-First Search (BFS)}

    \textbf{3. Stop when queue empty}

\end{frame}

% ------------------

\begin{frame}{Group walkthrough example}

    % We assume the source (obviously) here is A, with us mapping
    % all the routes from A to every other node existing in our tree.

    \includegraphics<1>[width=0.8\linewidth]{BFS1.jpg}
    \includegraphics<2>[width=0.8\linewidth]{BFS2.jpg}
    \includegraphics<3>[width=0.8\linewidth]{BFS3.jpg}
    \includegraphics<4>[width=0.8\linewidth]{BFS4.jpg}
    \includegraphics<5>[width=0.8\linewidth]{BFS5.jpg}
    \includegraphics<6>[width=0.8\linewidth]{BFS6.jpg}
    \includegraphics<7>[width=0.8\linewidth]{BFS7.jpg}
    \includegraphics<8>[width=0.8\linewidth]{BFS8.jpg}

\end{frame}

% ------------------

\begin{frame}{Why BFS Is Not Enough}

    \textbf{Edge Case: Weighted Graph}
    
        BFS assumes all edges have equal cost.
    
    \vspace{0.5em}
    
    Consider:
    
        \begin{itemize}
            \item $A \rightarrow B$ with weight 100
            \item $A \rightarrow C$ with weight 1
            \item $C \rightarrow B$ with weight 1
        \end{itemize}
    
    \vspace{0.5em}
    
    \textbf{What BFS does:}
        \begin{itemize}
            \item Finds path $A \rightarrow B$ (1 edge)
        \end{itemize}

\end{frame}

\begin{frame}{Why BFS Is Not Enough}

    \textbf{Actual shortest path:}
    \begin{itemize}
        \item $A \rightarrow C \rightarrow B$ with total cost $2$
    \end{itemize}
    
    \textbf{Key Limitation}
    
    \vspace{0.5em}
    
        BFS minimises the \textit{number of edges},  
        not the \textit{total path weight}.
    
    \vspace{0.8em}
    
    \textbf{Implication}
    
        When edge weights differ, BFS may return a suboptimal path.
    
    \vspace{0.8em}
    
    \textbf{Takeaway Point}
    
        Do not use BFS on graphs with weights on them!
    
\end{frame}

% ------------------

\begin{frame}{Priority Queues}
    
    \textbf{Motivation}
    
        When edge weights differ, we must expand the \textit{closest} node first.
    
    \vspace{0.8em}
    
    \textbf{Key Idea}
    
        Instead of processing nodes in FIFO order (like BFS),  
        we process nodes in order of \textit{smallest current distance}.
    
    \vspace{0.8em}
    
    \textbf{Data Structure}
    
        A \texttt{min-heap} (priority queue) allows us to:
        \begin{itemize}
            \item Always extract the node with smallest distance
            \item Efficiently update distances
        \end{itemize}
    
\end{frame}

% ------------------------------------------------

\section{Dijkstra - An Exploration}

% ------------------------------------------------

\begin{frame}{Why Would We Ever Use This?}

    \textbf{Problem}
    
        BFS only guarantees shortest path when all edges have equal weight.
    
    \vspace{0.8em}
    
    \textbf{Issue}
    
        When edge weights differ:
        \begin{itemize}
            \item The path with fewer edges may not have minimum total cost
            \item We must prioritise smaller cumulative distance
        \end{itemize}
    
    \vspace{0.8em}
    
    \textbf{Solution}
    
        Use Dijkstra’s Algorithm:
        \begin{itemize}
            \item Always expand the node with smallest known distance
            \item Guarantees shortest paths for non-negative weights
        \end{itemize}

\end{frame}

% ------------------

\begin{frame}{Algorithmic Theory}

    \textbf{Core Idea}
    
        Always expand the node with the smallest known distance.

\end{frame}

\begin{frame}{Algorithmic Theory}

    \vspace{0.8em}
    
    \textbf{Algorithm}
    
    \textbf{1. Initialize}
    \begin{itemize}
        \item Set $distance[v] = \infty$ for all nodes
        \item Set $distance[start] = 0$
        \item Create empty \texttt{min-heap} (priority queue)
        \item Insert $(0, start)$ into queue
    \end{itemize}
    
\end{frame}

\begin{frame}{Algorithmic Theory}

    \textbf{2. While queue not empty}
        \begin{itemize}
            \item Extract node $u$ with smallest distance
            \item If extracted distance is outdated, skip
            \item For each neighbour $v$ with edge weight $w$:
            \begin{itemize}
                \item If $distance[u] + w < distance[v]$:
                \begin{itemize}
                    \item Update $distance[v]$
                    \item Insert $(distance[v], v)$ into queue
                \end{itemize}
            \end{itemize}
        \end{itemize}
    
\end{frame}

\begin{frame}{Algorithmic Theory}

    \textbf{3. Stop when queue empty}

\end{frame}

\begin{frame}{Algorithmic Theory}

    \textbf{Time Complexity (Binary Min-Heap)}
    
        \begin{itemize}
            \item Best case: $O((V + E)\log V)$
            \item Average case: $O((V + E)\log V)$
            \item Worst case: $O((V + E)\log V)$
        \end{itemize}
    
    \vspace{0.6em}
    
    \textbf{Space Complexity}
    
        \begin{itemize}
            \item $O(V + E)$
        \end{itemize}
    
\end{frame}

% ------------------

\begin{frame}{Group walkthrough example}

\includegraphics<1>[width=0.8\linewidth]{Dijkstra1.jpg}
\includegraphics<2>[width=0.8\linewidth]{Dijkstra2.jpg}
\includegraphics<3>[width=0.8\linewidth]{Dijkstra3.jpg}
\includegraphics<4>[width=0.8\linewidth]{Dijkstra4.jpg}
\includegraphics<5>[width=0.8\linewidth]{Dijkstra5.jpg}
\includegraphics<6>[width=0.8\linewidth]{Dijkstra6.jpg}
\includegraphics<7>[width=0.8\linewidth]{Dijkstra7.jpg}
\includegraphics<8>[width=0.8\linewidth]{Dijkstra8.jpg}

\end{frame}

% ------------------

\begin{frame}{Edge Cases and Limitations}

    \textbf{1. Negative Edge Weights}
    
        \begin{itemize}
            \item Dijkstra assumes all edge weights are non-negative
            \item Negative weights can produce incorrect shortest paths
            \item Use Bellman–Ford instead
        \end{itemize}
    
    \vspace{0.6em}
    
    \textbf{2. Disconnected Graphs}
    
        \begin{itemize}
            \item Unreachable nodes remain at $distance = \infty$
        \end{itemize}

\end{frame}

\begin{frame}{Edge Cases and Limitations}

    \textbf{3. Duplicate Priority Queue Entries}

        \begin{itemize}
            \item A node may be inserted multiple times with different distances
            \item Outdated entries are skipped during extraction
        \end{itemize}
        
    \vspace{0.6em}
    
    \textbf{4. When to Use Another Algorithm}
    
        \begin{itemize}
            \item Equal weights → use BFS
            \item Weights restricted to $0$ or $1$ → use 0–1 BFS
        \end{itemize}

\end{frame}

% ------------------------------------------------

\section{0--1 BFS - An Exploration}

% ------------------------------------------------

\begin{frame}{Why Would We Ever Use 0--1 BFS?}

    \textbf{Observation}
    
        Dijkstra works for all non-negative weights,  
        but requires a priority queue and incurs a $\log V$ factor.

    \vspace{0.8em}
    
    \textbf{Special Case}
    
        If edge weights are restricted to $\{0,1\}$:
        \begin{itemize}
            \item We do not need full priority ordering
            \item Only two distance transitions are possible
        \end{itemize}
    
\end{frame}
    
\begin{frame}{Why Would We Ever Use 0--1 BFS?}
    
    \textbf{Key Idea}
    
        Exploit the weight structure to replace the heap  
        with a deque.
        
    \vspace{0.8em}
    
    \textbf{Guarantee}
    
        0--1 BFS computes shortest paths  
        when all edge weights are exactly $0$ or $1$.

\end{frame}

% ------------------

\begin{frame}{Algorithmic Theory}

    \textbf{Core Idea}
    
        Simulate Dijkstra without a heap by using a deque.
    
        \vspace{0.8em}
    
        Maintain $distance[v]$ as usual.

\end{frame}

\begin{frame}{Algorithmic Theory}
    \textbf{Algorithm}
    
        \textbf{1. Initialize}
            \begin{itemize}
                \item Set $distance[v] = \infty$ for all nodes
                \item Set $distance[start] = 0$
                \item Create empty \texttt{deque}
                \item Push $start$ to front
            \end{itemize}

\end{frame}

\begin{frame}{Algorithmic Theory}

    \textbf{2. While deque not empty}
        \begin{itemize}
            \item Pop node $u$ from front
            \item For each neighbour $v$ with weight $w \in \{0,1\}$:
            \begin{itemize}
                \item If $distance[u] + w < distance[v]$:
                \begin{itemize}
                    \item Update $distance[v]$
                    \item If $w = 0$: push $v$ to front
                    \item If $w = 1$: push $v$ to back
                \end{itemize}
            \end{itemize}
        \end{itemize}

\end{frame}

\begin{frame}{Algorithmic Theory}

    \textbf{3. Stop when deque empty}

\end{frame}

\begin{frame}{Algorithmic Theory}

    \textbf{Why It Works}
    
        The deque preserves nondecreasing distance order.
    
    \vspace{0.8em}
    
    \textbf{Time Complexity}
    
        \begin{itemize}
            \item Best case: $O(V + E)$
            \item Average case: $O(V + E)$
            \item Worst case: $O(V + E)$
        \end{itemize}
    
    \vspace{0.6em}
    
        \textbf{Space Complexity}
    
            \begin{itemize}
                \item $O(V + E)$
            \end{itemize}

\end{frame}


% ------------------

\begin{frame}{Group walkthrough example}

    \includegraphics<1>[width=0.8\linewidth]{01BFS1.jpg}

\end{frame}

% ------------------

\begin{frame}{Edge Cases and Limitations}

    \textbf{1. Weight Restriction}
    
        \begin{itemize}
            \item Only valid when all edge weights are exactly $0$ or $1$
            \item If weights include $2$ or higher, use Dijkstra (or Dial’s algorithm)
        \end{itemize}
    
    \vspace{0.6em}
    
        \textbf{2. Disconnected Graphs}
        
        \begin{itemize}
            \item Unreachable nodes remain at $distance = \infty$
        \end{itemize}

\end{frame}

\begin{frame}{Edge Cases and Limitations}
    
    \textbf{3. Multiple Relaxations}
    
        \begin{itemize}
            \item Nodes may be relaxed multiple times
            \item Do not permanently mark nodes as “visited” after first processing
        \end{itemize}
    
    \vspace{0.6em}
    
    \textbf{4. Negative Weights}
    
        \begin{itemize}
            \item Not valid if negative edges are present
        \end{itemize}

\end{frame}

% ------------------------------------------------

\section{Choosing the Right Algorithm (Quick Rules)}

% ------------------------------------------------

\begin{frame}{Choosing the Right Algorithm}

    \begin{itemize}
    
        \item \textbf{DFS:} Use when exploring graph structure such as connectivity or cycles, where shortest path distance is not required.
        
    \pause
        
        \item \textbf{BFS:} Use when all edges have equal weight and you need the shortest path in terms of number of edges.
    
    \pause
        
        \item \textbf{0--1 BFS:} Use when edge weights are restricted to exactly 0 or 1 and you need shortest path distances efficiently.
    
    \pause
        
        \item \textbf{Dijkstra:} Use when edge weights are non-negative and may vary, and you need the shortest path from a source node.
        
    \end{itemize}

\end{frame}

% ------------------------------------------------

\section{Problem Sheet}

% ------------------------------------------------


\begin{frame}{Problem Sheet}

    \begin{itemize}
    
        \item Solve the problems using the appropriate shortest-path algorithm
    
    \pause
    
        \item Work individually or in small groups
    
    \pause
    
        \item Refer to previous slides if needed
    
    \pause
    
        \item We will review selected solutions afterwards
    
    \end{itemize}

\end{frame}

% ------------------------------------------------

\section{Final Remarks}

% ------------------------------------------------

\begin{frame}{Further Reading}

    \begin{itemize}
    
        \item A document containing relevant supplementary + research reading will be available on the KCPC website (for this workshop) and the KCPC Discord Resource section for you to check out :D
    
    \end{itemize}
    
\end{frame}

% ------------------------------------------------

\end{document}